\section{Unipotent representations for finite groups of Lie type}\label{sec:finite_groups}

In this first section, we recall the necessary structural and representation theoretic properties of finite groups of Lie type leading to Lusztig's construction of the nonabelian Fourier transform described in the introduction. To this aim, suppose that $\mathbb{G}$ is a connected reductive group over $\overline{\FF}_p$ with Frobenius map $\Fr:\mathbb{G}\rightarrow\mathbb{G}$. The group $\GFr$ is a finite group of Lie type and we are interested in studying the set of irreducible unipotent representations of $\GFr$ denoted by $\Irr_{\un}(\GFr)$. We say that $\GFr$ is \textit{split} if the induced Frobenius action on the set of simple roots of $\GG$ is trivial. We shall assume throughout that $\GFr$ is a split group. We recall that Deligne--Lusztig induction associates, to each pair $(\TT,\theta)$ where $\TT$ is a $\Fr$-stable maximal tori of $\GG$ and $\theta$ is a character of $\TT^{\Fr}$, a virtual character $R_{\TT,\theta}$. Then we define 
$$\Irr_{\un}(\GFr)=\{\rho\in\Irr(\GFr)\ |\ \langle\rho,R_{\TT,\mathbf{1}}\rangle\neq0\text{ for some $\Fr$-stable torus $\TT$ of $\GG$}\}.$$
If the $\Fr$-maximal torus $\TT$ is maximally split and therefore contained in some $\Fr$-stable Borel subgroup $\BB$, then $R_{T_e,\mathbf{1}}=\Ind_{\BB^{\Fr}}^{\GG^{\Fr}}\mathbf{1}$, so the set of unipotent characters contain the principal series representations of $\GFr$. \textcolor{red}{By results of Lusztig, the classification of $\Irr_{\un}(\GFr)$ is reduced to the case when $\GG$ is adjoint simple.}

Let $W=N_{\GG}(\TT)/\TT$ be the Weyl group associated to $\GG$. \textcolor{red}{When the group $\GFr$ is split,} irreducible representations $\phi$ of $W$ are known to to parametrize principal series representations $\chi_\phi$ of $\GFr$, and we shall make use of this bijection throughout. \textcolor{red}{maybe mention here that other representation series are also in bijection with weyl groups of smaller types?}. In addition, for each $\phi\in\widehat{W}$, define the \textit{almost character} 
$$R_\phi:=\frac{1}{|W|}\sum_{w\in W}\phi(w)R_{\mathbb{T}_w,1},$$
arising as traces of certain character sheaves on the group $\GG$.
One can then partition $\Irr_{\un}(\GFr)$ into families by defining two representations $\rho_1,\rho_2$ to be related if 
$$\langle\rho_1,R_\phi\rangle\neq0\quad\text{and}\quad\langle\rho_2,R_\phi\rangle\neq0\quad\text{for some }\phi\in\widehat{W},$$
and then extending by transitivity. To each family $\mathcal{U}\subset\Irr_{\un}(\GFr)$, Lusztig associated a finite group $\Gamma_\mathcal{U}$, scalars $\Delta(x,\sigma)\in\{\pm1\}, (x,\sigma)\in\mathcal{M}(\Gamma_\mathcal{U})$ defined in \cite{Lusztig1984} and a bijection
\begin{align*}
    \mathcal{U}&\longleftrightarrow\mathcal{M}(\Gamma_\mathcal{U}):=\{(x,\sigma)\ |\ x\in\Gamma_\mathcal{U},\sigma\in\Irr(Z_{\Gamma_{\mathcal{U}}}(x))\}/\Gamma_\mathcal{U}\\
    \rho_{(x,\sigma)}&\longleftrightarrow(x,\sigma).
\end{align*}
The indexing set $\mathcal{M}(\Gamma_\mathcal{U})$ carries a natural pairing 
$$\left\{(x,\sigma),(y,\tau)\right\}=\frac{1}{|Z_{\Gamma_\mathcal{U}}(x)||Z_{\Gamma_\mathcal{U}}(y)|}\sum_{\substack{g\in\Gamma_{\mathcal{U}}\\xgyg^{-1}=gygx^{-1}}}\sigma(gyg^{-1})\tau(g^{-1}x^{-1}g),$$
and the above bijection satisfies the property that 
\begin{equation}\label{eqn:inner_product}
    \langle\rho_{(x,\rho)},R_\phi\rangle_{\GFr}=\begin{cases}
        \Delta(x,\sigma)\left\{(x,\rho),(y,\tau)\right\}&\text{ if }\chi_\phi=\rho_{(y,\tau)}\in\mathcal{U},\\
        0 & \text{ if }\chi_\phi\not\in\mathcal{U}.
    \end{cases}
\end{equation}
We recall that if $W$ is an irreducible Weyl group, then $\Delta(x,\rho)=1$ except for the families containing $\phi_{512,11}$ in $E_7$, $\phi_{4096,11}$ in $E_8$ or $\phi_{4096,26}$ in $E_8$. For any family $\mathcal{U}$, $\rho_{(e,1)}=\chi_\phi$ is a principal series representation corresponding to some representation $\phi$ of $W$. Such representations of $W$ are called \textit{special}, and they are the only ones satisfying the property that $\langle\chi,R_\phi\rangle\geq0$ for all unipotent characters $\chi$ (see \cite{Carter1985}*{\S12.3}). We define the \textit{unipotent almost characters} of $\GFr$ to be the set of orthonormal class functions
$$R_{(x,\sigma)}:=\sum_{(y,\tau)}\Delta(y,\tau)\{(y,\tau),(x,\sigma)\}\rho_{(y,\tau)},\quad(x,\sigma)\in\mathcal{M}(\Gamma_\mathcal{U}).$$
In particular, if $\chi_\phi=\rho_{(x,\rho)}$, then $R_\phi=R_{(x,\rho)}$. Lusztig also defined the map 
\begin{align*}
    \FT_{\Gamma_\mathcal{U}}:\CC[\mathcal{U}]&\longrightarrow\CC[\mathcal{U}]\\
    \rho_{(x,\sigma)}&\longmapsto\sum_{(y,\tau)\in\mathcal{M}(\Gamma_\mathcal{U})}\{(x,\sigma),(y,\tau)\}\rho_{(y,\tau)}
\end{align*}
The \textit{unipotent nonabelian Fourier transform} of $\GFr$ on the space $R_{\un}(\GFr)=\oplus_{\mathcal{U}\subseteq\Irr_{\un}(\GFr)}\CC[\mathcal{U}]$ is the map
$$\FT^{\GFr}:=\bigoplus_{\mathcal{U}\subset\Irr_{\un}(\GFr)}\FT_{\Gamma_\mathcal{U}},$$
and it satisfies many interesting properties. Most importantly, it is a self-Hermitian involution with respect to the standard inner product of characters and, in addition, $\FT^{\GFr}(\chi_\phi)=R_\phi$. More generally, it takes unipotent representations to the almost characters arising as traces of character sheaves on $\mathbb{G}$. Let us discuss some important examples that will be used in our tables in Section \ref{sec:tables}.

\begin{example}
    When $\GG=\PGL_n(\overline{\FF_p})$ and $\Fr$ is the standard $q$-th power Frobenius, then $\GFr=\PGL_n(\FF_q)$ and $W\cong S_n$. In this case, one can show that $R_{\phi}=\chi_{\phi}$ for all $\phi\in\widehat{S_n}$ so, in particular, all families have a unique element and unipotent representations coincide with the principal series. Thus, unipotent representations of $\GFr$ are naturally by partitions of $n$, and we shall use this notation in the tables from Section \ref{sec:tables}.
\end{example}

\begin{example}
    Suppose $\GFr=\mathrm{SO}_{2n+1}(\FF_q)$ or $\GFr=\mathrm{PGSp}_{2n}(\FF_q)$ (split, types $B_n$ or $C_n$), and with Weyl group $W\cong S_n\rtimes C_2^n$. Its representations are parametrized by double partitions of $n$, and we shall use this notation to label principal series representations. In turn, double partitions of $n$ are in bijection with symbols 
    \begin{equation*}
        \begin{pmatrix}
            \lambda_1,\ \ \lambda_2,\ \ \lambda_3, \ldots,\ \  \lambda_{k+s}\\
            \mu_1,\ \ \mu_2,\ \ \ldots,\mu_k\ \ 
        \end{pmatrix}
    \end{equation*}
    with $0\leq\lambda_1<\lambda_2<\cdots<\lambda_{k+s}$, $0\leq\mu_1<\mu_2<\cdots<\mu_k$, with $\lambda_1\mu_1\neq0$, rank $n$ \textcolor{red}{(should define rank or refer to a definition)} and defect $s=1$. To account for the remaining unipotent representations of $\GFr$, we allow for symbols with positive odd defect $s$. Remarkably, Lusztig showed that two unipotent representations lie in the same family if and only if the unordered set $\{\lambda_1,\ldots,\lambda_{k+s},\mu_1,\dots,\mu_k\}$ coincide. See \cite{Carter1985}*{\S11.4 and \S13.8} for a complete overview of the theory.

    For example, there are $10$ double partitions of $n=3$ (see Table \ref{tbl:F4_K4_1}) and the two symbols $\theta_1=\begin{psmallmatrix}
            0 \ 1\ 2\ 3\\
            \ \ \ 1\ \ \
        \end{psmallmatrix}$ and $\theta_2=\begin{psmallmatrix}
            0 \ 1\ 3\\
            \ \ - \ \
        \end{psmallmatrix}$ are the only two rank-$3$ symbols with higher defect (see Table \ref{tbl:F4_K4_2}).
        The $12$ unipotent representations are partitioned in $4$ trivial families and $2$ families of size $4$ containing the symbols
        \begin{align*}
            \left\{[1,11]=\begin{psmallmatrix}
                0 \ 1\ 3\\
                \ 1\ 2\
            \end{psmallmatrix},\ [\cdot,12]=\begin{psmallmatrix}
                0 \ 1\ 2\\
                \ 1\ 3\
            \end{psmallmatrix},\ [1^3,\cdot]=\begin{psmallmatrix}
                1 \ 2\ 3\\
                \ 0\ 1\
            \end{psmallmatrix},\ \theta_1=\begin{psmallmatrix}
                0 \ 1\ 2\ 3\\
                \ \ \ 1\ \ \
            \end{psmallmatrix}\right\}&
            \quad\text{and}\quad\\
            \left\{[2,1]=\begin{psmallmatrix}
                0 \ 3\\
                \ 1\
            \end{psmallmatrix},\ [\cdot,3]=\begin{psmallmatrix}
                0 \ 1\\
                \ 3\
            \end{psmallmatrix},\ [12,\cdot]=\begin{psmallmatrix}
                1 \ 3\\
                \ 0\
            \end{psmallmatrix},\ \theta_2=\begin{psmallmatrix}
                0 \ 1\ 3\\
                \ \ - \ \
            \end{psmallmatrix}\right\}&
        \end{align*}
        When $n=4$, there are $20$ double partitions (see Table \ref{tbl:F4_K1_1}) and $5$ rank-$4$ symbols with higher defect, namely $\theta_1=\begin{psmallmatrix}
            0 \ 1\ 2\ 3\ 4\\
            \ \ \ 1\ 2 \ \ \
        \end{psmallmatrix}$, $\theta_2=\begin{psmallmatrix}
            0 \ 1\ 2\ 3\\
            \ \ \ 2 \ \ \
        \end{psmallmatrix}$, $\theta_3=\begin{psmallmatrix}
            0 \ 1\ 2\ 4\\
            \ \ \ 1 \ \ \
        \end{psmallmatrix}$, $\theta_4=\begin{psmallmatrix}
            0 \ 2\ 3\\
            \ \ - \ \
        \end{psmallmatrix}$ and $\theta_5=\begin{psmallmatrix}
            0 \ 1\ 4\\
            \ \ - \ \
        \end{psmallmatrix}$ (see Table \ref{tbl:F4_K1_2}). These unipotent representations all lie in a family of size $4$ whose associated group is $\Gamma\cong C_2$. As an example, the Fourier transform for these families is 
        $$\begin{pmatrix}
            1/2 & 1/2 & 1/2 & 1/2 \\
            1/2 & 1/2 & -1/2 & -1/2 \\
            1/2 & -1/2 & 1/2 & -1/2 \\
            1/2 & -1/2 & -1/2 & 1/2 
        \end{pmatrix}.$$
\end{example}

\begin{example}
    Consider the split finite group of Lie type $\GFr=F_4(\FF_q)$. The Weyl group $W$ is a Coxeter group of order $1152$ and $25$ irreducible representations labelled as $\phi_{d,k}$, where $d$ is the dimension and $k$ is its $b$-invariant. The outer automorphism of $W$ swapping long and short roots acts on $\widehat{W}$ while preserving these invariants, so for some pairs $d,k$ one has two inequivalent representations labelled by $\phi_{d,k}'$ and $\phi_{d,k}''$. In addition to the principal series representations, $\Irr_{\un}(\GFr)$ contains $7$ cuspidal representations labelled by $F_4^I[1],F_4^{II}[1],F_4[-1],F_4[\theta],F_4[\theta^2],F_4[i]$ and $F_4[-i]$. Finally, it contains $5$ more unipotent representations arising as irreducible components of the parabolic induction of the unique unipotent cuspidal representation of a Levi of type $B_2$, labelled as $B_{2,\mathbf{1}}$, $B_{2,\varepsilon}$, $B_{2,\varepsilon'}$, $B_{2,\varepsilon''}$ and $B_{2,r}$.

    The set $\Irr_{\un}(\GFr)$ is then partitioned in $11$ families: there are $8$ families consisting of a single representation, $2$ families $\{\phi_{4,1},\phi_{2,4}',\phi_{2,4}'',B_2,\mathbf{1}\}$, $\{\phi_{4,13},\phi_{2,16}',\phi_{2,16}'',B_{2,\varepsilon}\}$ whose associated group is $\Gamma\cong C_2$ and labels $(1,1),(1,\varepsilon),(g_2,1),(g_2,\varepsilon)$, and one large family of size $21$ whose associated group is $\Gamma\cong S_4$. See Table \ref{tbl:F4family} for the representations and their labels in $\mathcal{M}(S_4)$.
    \begin{table}[h!]
        \centering
        \begin{adjustbox}{max width=\linewidth}
            \begin{tabular}{ccccccccccccccccccccc}
                $\phi_{12,4}$ & $\phi_{9,6}''$ & $\phi_{9,6}'$ & $\phi_{1,12}''$ & $\phi_{1,12}'$ & $\phi_{4,7}''$ & $\phi_{4,7}'$ & $\phi_{4,8}$ & $\phi_{6,6}'$ & $\phi_{6,6}''$ & $\phi_{16,5}$ & $B_{2,r}$ & $B_{2,\varepsilon'}$ & $B_{2,\varepsilon''}$ & $F_4[\theta]$ & $F_4[\theta^2]$ & $F_4[i]$ & $F_4[-i]$ & $F_4^I[1]$ & $F_4^{II}[1]$ & $F_4[-1]$ \\\hline 
                $(1,1)$ & $(g_2',1)$ & $(1,\lambda^1)$ & $(g_2',\varepsilon')$ & $(1,\lambda^2)$ & $(g_4,1)$ & $(g_2,\varepsilon'')$ & $(g_2',\varepsilon'')$ & $(g_3,1)$ & $(1,\sigma)$ & $(g_2,1)$ & $(g_2',r)$ & $(g_4,-1)$ & $(g_2,\varepsilon')$ & $(g_3,\theta)$ & $(g_3,\theta^2)$ & $(g_4,i)$ & $(g_4,-i)$ & $(g_2',\varepsilon)$ & $(1,\lambda^3)$ & $(g_2,\varepsilon)$\\                                
            \end{tabular}
        \end{adjustbox}
        \caption{Largest $F_4(\FF_q)$-family of unipotent representations and their labels}
        \label{tbl:F4family}
    \end{table}
\end{example}


\newpage

When $\GFr$ is split (as it is the case throughout this document) it is convenient to express Lusztig's nonabelian Fourier transform in term of stable combinations of unipotent representations. This makes the verification of Conjecture \ref{conj:main} cleaner and more straightforward. Let $\mathcal{U}$ be a family of unipotent representations and for each $x\in\Gamma_\mathcal{U}$, define the virtual combinations of unipotent characters
\begin{equation}\label{eqn:virtual1}
    \Pi_\mathcal{U}(x,y)=\sum_{\sigma\in\widehat{Z_{\Gamma_\mathcal{U}}(x)}}\sigma(y^{-1})\rho_{(x,\sigma)},\quad\text{where } y\in Z_{\Gamma_\mathcal{U}}(x).
\end{equation}
Dually, when $\Gamma_\mathcal{U}$ is abelian, define the virtual combinations
\begin{equation}\label{eqn:virtual2}
    \Pi_\mathcal{U}(\sigma,\tau)=\sum_{y\in\Gamma_\mathcal{U}}\tau(y)\rho_{(y,\sigma)},\quad\text{if } \sigma,\tau\in \widehat{\Gamma_\mathcal{U}}.
\end{equation}


\begin{lemma}[\cite{AubertCiubotaruRomano2024}*{Lemma 5.1}]
    With the notation of \eqref{eqn:virtual1} and \eqref{eqn:virtual2}, 
    \begin{equation*}
        \FT_{\Gamma_\mathcal{U}}(\Pi_\mathcal{U}(x,y))=\Delta(x,y)\Pi_{\mathcal{U}}(y,x),\quad \FT_{\Gamma_\mathcal{U}}(\Pi_\mathcal{U}(\sigma,\tau))=\Pi_{\mathcal{U}}(\tau,\sigma),
    \end{equation*}
    the latter when $\Gamma_\mathcal{U}$ is abelian, where $\Delta(x,y)=1$ unless $\mathcal{U}$ is one of the families described in \eqref{eqn:inner_product}, in which case $\Delta(1,g_2)=\Delta(g_2,1)=-1$ and $1$ otherwise.
\end{lemma}


